\documentclass[aspectratio=169,notheorems]{corvinusmetropolis}
\newtheorem{theorem}{{Tétel.}}
\newtheorem{proposition}[theorem]{{Állítás.}}
\newtheorem{lemma}[theorem]{{Lemma.}}
\newtheorem{FA}[theorem]{{Az algebra alaptétele.}}
\newtheorem{corollary}[theorem]{{Következmény.}}
\theoremstyle{definition}
\newtheorem{definition}[theorem]{{Definíció.}}
\newtheorem{defprop}[theorem]{{Definíció-Állítás.}}
\newtheorem{notes}[theorem]{{Megjegyzés.}}

\email{magyarkuti@uni-corvinus.hu}
\author[Magyarkuti]{Gyula~Magyarkuti\,\orcidlink{0000-0003-4444-0481}}
\institute[Corvinus Egyetem]%
{%\inst{1}
  Budapesti Corvinus Egyetem\\
  Matematika Tanszék\\
  Sóház, 208/B

  \begin{minipage}[b]{0.8\linewidth}
  email: \putemail\\
  letölthető: \url{https://magyarkuti.github.io/algebra-1}
  \end{minipage}
  \hfill{\adjustbox{valign=B}{\qrcode[hyperlink]{https://magyarkuti.github.io/algebra-1}}}%% iff qrcode needed, otherwise comment this line!
}
\date[Tavasz]{\ontoday}%% overwrite later, if needed.

\titlegraphic{\hfill{\includegraphics[width=.2\linewidth]{\myinstlogo}}}

\renewcommand{\emph}{\alert}
\renewcommand{\mathbf}{\mathbb}
\newcommand{\lin}{\operatorname{lin}}
\newcommand{\gen}{\operatorname{gen}}
\newcommand{\im}{\operatorname{Im}}
\newcommand{\id}{\operatorname{id}}
\newcommand{\sgn}{\operatorname{sgn}}
\newcommand{\tr}{\operatorname{tr}}
\newcommand{\codim}{\operatorname{codim}}
%newcommand{\deg}{\operatorname{deg}}
\newcommand{\ip}[2]{\langle#1,#2\rangle}
\newcommand{\uj}{\text{új}}
\newcommand{\rgi}{\text{régi}}
\newcommand{\Star}[1]{#1\ensuremath{^*}\kern-\scriptspace}
\newcommand{\CStar}{\Star{\ensuremath{\mathrm{C}}}}
\endinput


\title[Algebra 1.]{Vektorrendszerek}
\subtitle{lineárisan összefüggősége, függetlensége}
\date[8. hét]{2025.~november 9-én}
\renewcommand{\emph}[1]{\alert{\textit{#1}}}
\begin{document}
\openingframes{Tartalom}

\section{Lineárisan független rendszerek}
\frame{\frametitle{lineárisan összefüggő vektorrendszer}
\metroset{block=fill}
\begin{definition}
	Egy véges $\left\{ y_1,\ldots,y_n \right\}$ vektorrendszert \emph{lineárisan összefüggőnek}
	mondunk, ha van olyan vektora, amely kifejezhető a többi vektor lineáris kombinációjaként.
\end{definition}
\metroset{block=transparent}
Úgy is fogalmazhatnánk, hogy az $\left\{ y_1,\ldots,y_n \right\}$ rendszer pontosan akkor
lineárisan összefüggő, ha létezik $1\leq k\leq n$ index, amelyre
\[
	y_k\in\lin\left\{ y_1,\ldots,y_{k-1},y_{k+1},\ldots,y_n \right\}.
\]
}
\frame{
Nyilvánvaló példa, hogy ha a vektorrendszer a $0$ vektort tartalmazza,
akkor lineárisan összefüggő.

Hasonlóan, ha ugyanaz a vektor többször szerepel a vektorrendszerben, a rendszer szintén összefüggő.

Olyan vektorrendszerre, amelyre nem igaz a lineárisan összefüggőség definíciója példa egy egyedüli nem zérus vektor,
de az $\left\{  \right\}$ üres rendszer is.
}
\frame{
\begin{proposition}
	Legyen $\left\{ y_1,\ldots,y_n \right\}$ vektorrendszer rögzítve a $V$ vektortérben.
	A vektorrendszerre tett alábbi feltevések egymással ekvivalensek.
	\begin{enumerate}
		\item Lineárisan összefüggő;
        \item Van olyan elem a vektorrendszer lineáris burkában, 
            amely nem csak egyféleképpen áll elő mint az $y_1,\ldots,y_n$
		      vektorok lineáris kombinációja, azaz:\\
              létezik 
              \(
              z\in \lin\left( \left\{ y_1,\ldots,y_n \right\} \right), 
              \)
              amelyre 
              \[
                  z=\sum_{j=1}^n\xi_jy_j \quad\text{ és }\quad
                  z=\sum_{j=1}^n\eta_jy_j
              \]
		      és létezik $1\leq k\leq n$, amelyre $\xi_k\neq\eta_k$.
		\item Vannak olyan nem mind zérus $\alpha_1,\ldots,\alpha_n$ skalárok, amelyekkel
		      \(
			      \sum_{j=1}^n\alpha_jy_j=0,
		      \)
              azaz:\\
              a vektorrendszernek van nem triviális lineáris kombinációja, amely a zérus vektort eredményezi.
              \qedhere
	\end{enumerate}
\end{proposition}
%\begin{proof}
%	Körben bizonyítunk.
%	\begin{description}
%		\item[$1.\Rightarrow 2.$]
%		      Tegyük fel, hogy $y_k=\sum_{j=1}^{k-1}\eta_jy_j+\sum_{j=k+1}^n\eta_jy_j$.
%		      Ekkor az alábbi együttható rendszerek
%		      \[
%			      \left( \eta_1,\ldots,\eta_{k-1},0,\eta_{k+1},\ldots,\eta_n \right)
%			      \qquad
%			      \left( 0,\ldots,0,1,0,\ldots,0 \right)
%		      \]
%		      a $k$-adik helyen biztosan különböznek,
%		      hiszen $0\neq 1$,
%		      és mind a két együttható rendszerrel képzett lineáris kombináció ugyanazt az $y_k$ vektort eredményezi.
%		\item[$2.\Rightarrow 3.$]
%		      Világos, hogy
%		      \[
%			      0=z-z=
%			      \sum_{j=1}^n\left( \xi_j-\eta_j \right)y_j
%		      \]
%		      és a $k$-adik skalár nem zérus.
%		\item[$3.\Rightarrow 1.$]
%		      Tegyük fel most, hogy
%		      \(
%		      \sum_{j=1}^n\alpha_jy_j=0,
%		      \)
%		      és, hogy $\alpha_k\neq 0.$
%		      Ekkor
%		      \[
%			      y_k=\sum_{\substack{=1\\j\neq k}}^n-\frac{1}{\alpha_k}\alpha_jy_j
%		      \]
%		      azaz a $k$-adik vektor tekinthető mint a többi vektor valamely lineáris kombinációja.
%	\end{description}
%	Ezt kellett belátni.
%\end{proof}
}
\frame{
Fontos észrevétel a következő.
\begin{proposition}
	Minden lineárisan összefüggő vektorrendszert tartalmazó vektorrendszer maga is lineárisan összefüggő.
\end{proposition}

\begin{definition}
    Egy nem véges vektorrendszert \emph{lineárisan összefüggőnek} nevezünk,
	ha van véges részrendszere, amely lineárisan összefüggő.
\end{definition}
Persze továbbra is igaz, hogy egy összefüggő rendszert tartalmazó rendszer is összefüggő marad.

A lineárisan összefüggőség a lineáris burok operációval is megfogható:
\begin{proposition}
    Egy vektortér $H$ vektorrendszere pontosan akkor lineárisan összefüggő, ha létezik $y\in H$, amelyre $y\in\lin(H\smallsetminus\left\{ y \right\})$.
\end{proposition}
}
%\begin{proof}
%    Ha $H$ lineárisan összefüggő, akkor van $\left\{ x_1,\ldots,x_k \right\}$ véges összefüggő vektorrendszere.
%    Itt az egyik vektor előáll mint a többi lineáris kombinációja. 
%    Ha ezt a vektort $y$ jelöli, akkor $y\in\lin\left( H\smallsetminus\left\{ y \right\} \right)$.
%    Megfordítva, ha valamely $y\in H$ mellett van $\left\{ x_1,\ldots,x_k \right\}\subseteq H\smallsetminus\left\{ y \right\}$ vektorrendszer, amelynek egy lineáris kombinációja éppen $y$, 
%    akkor az $\left\{ x_1,\ldots,x_k,y \right\}$ véges rendszer egy lineárisan összefüggő rendszere $H$-nak, ergo $H$ valóban lineárisan összefüggő.
%\end{proof}
\frame{\frametitle{lineárisan független vektorrendszer}
\metroset{block=fill}
\begin{definition}
	Egy vektorrendszer \emph{lineárisan független}, ha nem lineárisan összefüggő.
\end{definition}
\metroset{block=transparent}
Az előző állítás tagadásával azonnal meg is kapjuk a lineárisan független rendszerek jellemzését
a lineáris burok operáció segítségével.
\begin{proposition}
    Egy vektortér $H$ vektorrendszere pontosan akkor lineárisan független, ha minden $y\in H$ mellett $y\notin\lin(H\smallsetminus\left\{ y \right\})$ is teljesül.
\end{proposition}
Fontos észrevétel a következő:
\begin{proposition}
    Egy lineárisan független rendszer minden részrendszere is lineáris független marad.
\end{proposition}
}
\frame{
A definíció szerint egy nem véges vektorrendszer akkor lineárisan független, ha minden véges részrendszere is az.
Egy véges vektorrendszer lineárisan függetlenségét, pedig a következő egymással ekvivalens állítások karakterizálják.
\begin{proposition}
	Legyen $\left\{ y_1,\ldots,y_n \right\}$ vektorrendszer rögzítve a $V$ vektortérben.
	A vektorrendszerre tett alábbi feltevések egymással ekvivalensek.
	\begin{enumerate}
		\item Lineárisan független;
		\item A vektorrendszer lineáris burkában minden elem csak egyféleképpen áll elő,
		      mint az $y_1,\ldots,y_n$ vektorok lineáris kombinációja.
		\item Az $y_1,\ldots,y_n$ vektoroknak csak a \emph{triviális lineáris kombinációja}\index{triviális lineáris kombináció} zérus,
		      azaz
		      \[
			      \sum_{j=1}^n\alpha_jy_j=0\text{ esetén }\alpha_1=\alpha_2=\dots=\alpha_n=0.\qedhere
		      \]
	\end{enumerate}
\end{proposition}
}
\frame{
Egy a zéró vektort tartalmazó vektorrendszer persze lineárisan összefüggő,
és egy nem zérus vektorból álló egyelemű vektorrendszer lineárisan független.

A következő állítás sokszor teszi kényelmessé a gondolatmenetünket.
\begin{proposition}
	Legyen $\left\{ y_1,\ldots,y_n \right\}$ egy olyan legalább két elemű vektorrendszer,
	amelynek első eleme nem a zérus vektor, tehát $y_1\neq 0$.
	A vektorrendszer pontosan akkor lineárisan összefüggő,
	ha létezik olyan eleme,
	amely pusztán az előző elemek lineárisan kombinációja.

	Formálisabban: akkor és csak akkor,
	ha létezik olyan $k$ index, amelyre 
	$2\leq k\leq n$ és $y_k\in\lin\left\{ y_1,\ldots,y_{k-1} \right\}$.
\end{proposition}
}
%\begin{proof}
%	Tegyük fel, hogy a vektorrendszer lineárisan összefüggő.
%	Ekkor van olyan a zérus vektort eredményező lineáris kombinációja
%	\(
%	\alpha_1y_1+\dots+\alpha_ny_n=0,
%	\)
%	ahol nem az összes együttható nulla.
%	Legyen $k$ a lineáris kombinációban a legnagyobb nem nulla együttható indexe.
%	Világos, hogy $k\neq 1$,
%	hiszen $y_1\neq 0$.
%	Persze a $k$ feletti együtthatók mind nullák,
%	emiatt
%	\[
%		\alpha_1y_1+\dots+\alpha_ky_k=0.
%	\]
%	Itt már $\alpha_k\neq 0$, tehát $y_k$ kifejezhető az előző vektorok segítségével:
%	\[
%		y_k=
%		\frac{-1}{\alpha_k}\alpha_1y_1+\dots+\frac{-1}{\alpha_{k-1}}\alpha_{k-1}y_{k-1}.\qedhere
%	\]
%\end{proof}
\section{Maximális és minimális független és összefüggő rendszerek}
\frame{
A lehető legszűkebb lineárisan független rendszer az $\left\{  \right\}$ üres vektorrendszer,
amelynek egyetlen eleme sincs.
A lehető legbővebb generátorrendszer az egész vektortér.
Ennél sokkal érdekesebb és fontosabb a lineárisan független rendszerek közül a lehető legbővebbet keresni,
és generátorrendszerek közül a lehető legszűkebbet keresni.
\begin{definition}\index{maximális lineárisan független rendszer}\index{minimális generátorrendszer}
	Egy lineárisan független rendszert \emph{maximális lineárisan független rendszernek} nevezünk,
	ha nem lehet bővíteni úgy, hogy lineárisan független maradjon.

	Egy generátorrendszert \emph{minimális generátorrendszernek} mondunk, ha nem lehet szűkíteni úgy,
	hogy generátorrendszer maradjon.
\end{definition}
}
\frame{
\begin{proposition}
    Legyen $H$ egy (nem feltétlen véges) vektorrendszer a V vektortérben.
    \begin{enumerate}
        \item Tegyük fel, hogy $y\notin H$ és $H$ lineárisan független.
            \\
            A $H\cup\left\{ y \right\}$ pontosan akkor lineárisan összefüggő, ha $y\in\lin H$.
        \item Tegyük fel, hogy $y\in H$ és $H$ generátorrendszer.
            \\
            A $H\smallsetminus\left\{ y \right\}$ pontosan akkor generátorrendszer, ha $y\in\lin\left( H\smallsetminus \left\{ y \right\} \right)$.
            \qedhere
    \end{enumerate}
\end{proposition}
%\begin{proof}
%    Voltaképpen négy különböző állítást kell igazolnunk:
%    \begin{description}
%        \item[$1.\Leftarrow$]
%            Ha $y\in\lin H$, akkor $H$-nak van $\left\{ x_1,\ldots,x_k \right\}$ véges részrendszere,
%            amelyre az $\left\{ x_1,\ldots,x_k,y \right\}$ vektorrendszer $y$ elemére igaz, hogy a többi elem lineáris kombinációja,
%            ergo $H\cup\left\{ y \right\}$ lineárisan összefüggő.
%        \item[$1.\Rightarrow$]
%            Ha $H\cup\left\{ y \right\}$ lineárisan összefüggő, akkor van $H\cup\left\{ y \right\}$-nak véges részrendszere, amely lineárisan összefüggő.
%            Ebben a rendszerben $y$-nak szerepelnie kell hiszen $H$ lineárisan független.
%            Azt kapjuk tehát, hogy bizonyos $x_1,\ldots,x_k\in H$ vektorok mellett az $\left\{ x_1,\ldots,x_k,y \right\}$ vektorrendszer lineáris összefüggő.
%            No de, $H$ lineárisan független, ezért csak az utolsó vektor, az $y$ vektor lehet az előző vektorok lineáris burkában.
%        \item[$2.\Rightarrow$]
%            Ha $H\smallsetminus\left\{ y \right\}$ generátorrendszer, akkor minden vektor így az $y$ is eleme a $\lin\left( H\smallsetminus\left\{ y \right\} \right)$ altérnek.
%        \item[$2.\Leftarrow$]
%            Ha $y\in\lin\left( H\smallsetminus\left\{ y \right\} \right)$, akkor $H\subseteq\lin\left( H\smallsetminus\left\{ y \right\} \right)$.
%            Emiatt $\lin H\subseteq\lin\left( \lin\left( H\smallsetminus\left\{ y \right\} \right) \right)\subseteq\lin\left( H\smallsetminus\left\{ y \right\} \right)$.
%            Ebből persze $\lin\left( H\smallsetminus\left\{ y \right\} \right)=\lin H$ következik.
%            No de, $H$ generátorrendszer, ergo $H\smallsetminus\left\{ y \right\}$ is az.
%            \qedhere
%    \end{description}
%\end{proof}
A előző állítás 1. pontját úgy érdemes értelmeznünk, hogy a $V$ vektortér egy $H\subseteq V$ lineáris független rendszere akkor és csak akkor nem bővíthető a lineárisan függetlenség megtartásával, ha $\lin H=V$, azaz $H$ egy generátorrendszer.
\\
A 2. pont is hasonló: A $H$ generátorrendszer egyetlen eleme sem elhagyható a generátorrendszer tulajdonság megtartásával akkor és csak akkor, ha minden $y\in H$ mellett $y\notin\lin\left( H\smallsetminus\left\{y \right\} \right)$, azaz $H$ egy lineárisan független rendszer.
}
\frame{
Bebizonyítottuk tehát az alábbi fontos állítást.
\metroset{block=fill}
\begin{proposition}
	Egy vektortér egy vektorrendszerére az alábbi feltevések ekvivalensek.
	\begin{enumerate}
		\item A vektorrendszer maximális lineárisan független rendszer.
		\item A vektorrendszer egyszerre lineárisan független és generátorrendszer.
		\item A vektorrendszer minimális generátorrendszer.\qedhere
	\end{enumerate}
\end{proposition}
}
\metroset{block=transparent}
\frame{\frametitle{Független rendszer csere}
\begin{lemma}\label{le:fgtlncsere}
	Legyen $\left\{ y_1,\ldots,y_n \right\}$ egy lineárisan független rendszere valamely vektortérnek,
	és tegyük fel, hogy a tér valamely $x$ vektorára
	\[
		x=\sum_{j=1}^n\xi_jy_j,
	\]
	ahol $\xi_k\neq 0$ az egyik $1\leq k\leq n$ mellett.

	Ekkor az $x$ vektor becserélhető a $k$-adik helyen a független rendszerbe
	úgy, hogy az független maradjon, azaz az
	\[
		\left\{ y_1,\ldots,y_{k-1},x,y_{k+1},\ldots,y_n \right\}
	\]
	vektorrendszer is lineárisan független.
\end{lemma}
%\begin{proof}
%	Megmutatjuk, hogy a cserélt rendszernek egyedül a triviális lineáris kombinációja zérus.
%	Legyen tehát
%	\[
%		0
%		=\sum_{\substack{j=1\\j\neq k}}^n\eta_jy_j+\eta_kx
%		=\sum_{\substack{j=1\\j\neq k}}^n\eta_jy_j+\sum_{j=1}^n\eta_k\xi_jy_j
%		=\sum_{\substack{j=1\\j\neq k}}^n\left( \eta_j+\eta_k\xi_j \right)y_j+\eta_k\xi_ky_k.
%	\]
%	Mivel az eredeti rendszer lineárisan független,
%	ezért az utóbbi lineáris kombináció minden együtthatója a test null eleme.
%	A $k$-adikkal kezdve $\eta_k\xi_k=0$, $\xi_k\neq 0$, így $\eta_k=0$.
%	Ezt a nem $k$-adik együtthatókba visszahelyettesítve
%	már azt kapjuk, hogy $\eta_j=0$ minden $j\neq k$ mellett is.
%\end{proof}
}
\frame{}
\end{document}
